\documentclass[11pt]{article}
\usepackage[utf8]{inputenc}
\usepackage{amsmath}
\usepackage{amssymb}
\usepackage{geometry}
\usepackage{graphicx}
\usepackage{booktabs}
\usepackage{array}
\usepackage[hidelinks]{hyperref}
\geometry{a4paper, margin=1in}

\newcommand{\reff}{r_{\mathrm{eff}}}
\newcommand{\Keff}{K_{\mathrm{eff}}}
\newcommand{\rbase}{r_{\mathrm{base}}}

\title{Model-Based Offline RL for AM with Population Models}

\begin{document}

\maketitle

\setcounter{tocdepth}{2}
\tableofcontents
\newpage

\section{Problem Setting (The hidden parameter MDP Formulation)}
The problem is framed as a restricted Partially Observable Markov Decision Process (POMDP) where the state space is factored into a fully observable component and a completely hidden, stationary component.

\begin{itemize}
    \item \textbf{Factored State Space ($\mathcal{S} = \mathcal{X} \times \mathcal{M}$):}
    \begin{itemize}
        \item $\mathcal{X}$ is the set of fully observable, discrete abundance intervals (e.g., $x \in \{0, 1, \dots, 100\}$).
        \item $\mathcal{M}$ is the hidden model parameter space. In this case, $\mathcal{M}$ represents the specific intrinsic growth rate $r$ driving the current environment. $r$ is stationary for a given episode or species but remains hidden from the agent.
    \end{itemize}
    
    \item \textbf{Action Space ($\mathcal{A}$):} A finite set of discrete management actions (e.g., $a \in \{0, 1, 2, 3\}$). Each action is mapped to a specific carrying capacity $K(a)$.
    
    \item \textbf{Observations ($\Omega$):} Assuming zero observation noise ($\epsilon_t = 0$) on the discrete bins, the observation $o_t$ is exactly the discrete state $x_t$. The parameter $r$ remains unobserved.
    
    \item \textbf{Transition Dynamics ($\mathcal{T}$):} The probability of transitioning to the next discrete abundance state $x_{t+1}$ depends strictly on the current discrete state $x_t$, the action $a_t$, and the hidden model $r$: 
    $$P(x_{t+1} | x_t, a_t, r)$$
\end{itemize}

\noindent \textbf{Comparison to the hmMDP Framework:} Traditional hidden-model MDP (hmMDP) / mixed-observability solvers for adaptive management~\cite{chades2012momdp, memarzadeh2018} restrict $\mathcal{X}$ to very few states (e.g., 2 states) to analytically find a small set of models $\mathcal{M}$ that span all optimal policies. This formulation expands $\mathcal{X}$ to 50--1000 states. Because exact analytical solvers fail at this scale, analytical spanning is replaced with Model-Based Offline RL, which learns an empirical transition function directly from a dataset.

\section{The Simulation Environment (Data Generation)}
To generate the offline dataset $\mathcal{D} = \{(x_t, a_t, x_{t+1})\}$, the simulation environment runs in continuous space but outputs discrete observations.

\paragraph{Continuous Hidden Update}
$$s_{t+1} = s_t \exp\left(r \left(1 - \frac{s_t}{K(a_t)}\right)\right)$$
\textit{(Note: Environmental process noise $\eta_t = 0$. Open for further noise}

\paragraph{Discretization (The ``Observation'')}
The continuous abundance $s_{t+1}$ is mapped to a discrete state index $x_{t+1}$ using a defined set of intervals (e.g., $x_{t+1} = \lfloor s_{t+1} / \text{bin\_width} \rfloor$).

\paragraph{Data Logging}
The environment logs the tuple $(x_t, a_t, x_{t+1})$. The true continuous state $s_t$ and the intrinsic growth rate $r$ are discarded and remain hidden from the final dataset.

\section{The Model-Based Deep RL Pipeline}
The RL architecture operates across three distinct phases to ensure scalability, safety, and adaptability.

\subsection{Phase 1: Offline Model Learning (The Black-Box Ensemble)}
As the state space is too large for simple empirical counting, the methodology relies on deep function approximation.
\begin{itemize}
    \item The system is treated as a pure black box; no prior knowledge of the true parameters $r$, $K(a_t)$, or the underlying ecological equation (e.g., the Ricker model) is assumed.
    \item Instead of a discrete probability table, an ensemble of $N$ deep neural networks $f_{\theta_i}(x_t, a_t)$ is used to approximate the transition dynamics $\hat{P}(x_{t+1} | x_t, a_t)$.
    \item Each network in the ensemble is trained independently on bootstrapped subsets of the offline dataset $\mathcal{D}$ to capture epistemic uncertainty.
\end{itemize}

\subsection{Phase 2: Safe Planning via Pessimism (Handling OOD Shifts)}
When training data and the real environment differ (Distribution Shift), a standard RL agent may make confident but catastrophic mistakes. In ecology, an Out-of-Distribution (OOD) mistake can result in the collapse of a species. This risk is mitigated using uncertainty-penalized planning (e.g., the MOPO~\cite{yu2020mopo} or MOReL~\cite{kidambi2020morel} algorithms).

\begin{itemize}
    \item \textbf{Quantifying Epistemic Uncertainty:} If an action is considered in a well-represented state, all $N$ networks will interpolate safely and predict a similar next-state distribution. Conversely, in an OOD state, the networks will extrapolate in different directions, resulting in high variance.
    \item \textbf{Applying the Pessimism Penalty:} When the ensemble disagrees, a ``pessimistic'' reward penalty proportional to the predictive variance is applied. This ensures the planning algorithm avoids uncertain actions and restricts the agent to regions of the state-space well-represented in the offline data.
\end{itemize}

\subsection{Phase 3: Online Adaptation (Adaptive Management)}
Freezing the model ($\hat{\mathcal{T}}$) and policy ($\pi^*$) after the offline phase would defeat the purpose of \textbf{Adaptive Management (AM)}. The method requires the agent to update during execution to bridge the offline-to-online gap.

\begin{itemize}
    \item \textbf{Contrast with Multi-Model hmMDPs:} Traditional AM approaches define a discrete set of candidate models and update a belief state over them. This approach is computationally expensive and brittle if the true parameters fall outside the predefined set.
    \item \textbf{The Black-Box Continuous Update:} Instead of relying on predefined matrices, the ensemble of neural networks is maintained. As new real-world observations $(x_t, a_t, x_{t+1})$ are collected, online gradient updates are performed on the ensemble's weights. 
    \item \textbf{Closing the Loop:} As the models update with new data, predictions converge. The uncertainty variance drops, the pessimism penalty is lifted, and the agent can confidently manage newly discovered parts of the environment.
\end{itemize}

\newpage
% \section*{Part 1: The Environment Setup (Data Generator)}

% To generate moderate amounts of offline data, the intrinsic growth rate $r$ is fixed, and 4 distinct management actions are defined to alter the carrying capacity $K$.

% \subsection*{1. The Hidden Biological Parameters}
% \begin{itemize}
%     \item \textbf{Intrinsic Growth Rate ($r$):} The value is set to $r = 0.3$. This represents a moderately fast-growing species that recovers well but remains vulnerable to over-management.
%     \item \textbf{Max Population / Carrying Capacity ($K$):} The discrete action space $\mathcal{A} \in \{0, 1, 2, 3\}$ is mapped to specific carrying capacities:
%     \begin{itemize}
%         \item \textbf{Action 0 (Do Nothing):} $K(0) = 500$.
%         \item \textbf{Action 1 (Harvest/Degrade):} $K(1) = 300$.
%         \item \textbf{Action 2 (Moderate Support):} $K(2) = 700$.
%         \item \textbf{Action 3 (Aggressive Conservation):} $K(3) = 900$.
%     \end{itemize}
% \end{itemize}

% \subsection*{2. The State Space Discretization}
% \begin{itemize}
%     \item \textbf{Maximum Abundance:} 1000.
%     \item \textbf{Discrete States ($\mathcal{X}$):} 100 intervals of width 10.
% \end{itemize}

% \subsection*{3. The Trajectory Generation Loop}
% The environment executes the following loop to generate the offline data:
% \begin{enumerate}
%     \item A random continuous starting state (e.g., $s_0 = 450.5$) is sampled.
%     \item An action (e.g., $a_0 = 1$) is selected via a suboptimal policy.
%     \item The exact next continuous state is calculated using the Ricker growth equation.
%     \item The states are discretized for logging: $x_t = \lfloor s_t / 10 \rfloor$.
%     \item The tuple is logged into the dataset $\mathcal{D}$.
% \end{enumerate}

\phantomsection\addcontentsline{toc}{section}{Part 1: The Environment Setup (Data Generator)}
\section*{Part 1: The Environment Setup (Data Generator)}
In this redesigned setting, the agent's actions alter both the intrinsic growth rate ($r$) and the carrying capacity ($K$) via additive perturbations, which is ecologically more realistic than altering $K$ alone --- most real management interventions (predator control, supplementary feeding, captive breeding) act on vital rates rather than on habitat size. The base intrinsic growth rate $r_{\mathrm{base}}$ is treated as a \emph{hidden} parameter, fixed within an episode but resampled across episodes, formalising the Bayes-Adaptive (hmMDP) structure of the problem.

\phantomsection\addcontentsline{toc}{subsection}{1. Hidden and Known Biological Parameters}
\subsection*{1. Hidden and Known Biological Parameters}\begin{itemize}
    \item \textbf{Hidden Intrinsic Growth Rate ($r_{\mathrm{base}}$):} Drawn at the start of each episode from a known prior, $r_{\mathrm{base}} \sim \mathcal{U}(0.95, 1.00)$, then held fixed. The agent never observes $r_{\mathrm{base}}$ directly; it must infer it from observed transitions during planning.
    \item \textbf{Base Carrying Capacity ($K_{\mathrm{base}}$):} Known and fixed at $K_{\mathrm{base}} = 500$.
    \item \textbf{Effective parameters per timestep:} Each action $a_t$ contributes additive deltas $\Delta r(a_t)$ and $\Delta K(a_t)$, giving
    \[
        r_{\mathrm{eff}}(a_t) = r_{\mathrm{base}} + \Delta r(a_t), \qquad K_{\mathrm{eff}}(a_t) = K_{\mathrm{base}} + \Delta K(a_t).
    \]
\end{itemize}

\phantomsection\addcontentsline{toc}{subsection}{2. Action Space}
\subsection*{2. Action Space}
Each action specifies an ecological intervention with associated $(\Delta r, \Delta K)$ effects and a real-valued cost $c(a)$. Harvest actions carry \emph{negative} costs (i.e.\ revenue); conservation actions carry positive costs.

\subsubsection*{2.1 Smaller Configuration (5 Actions)}

The smaller set is used for initial pipeline validation. Each action exercises a different code path, covering every sign combination of $(\Delta r, \Delta K)$.

\begin{center}
\begin{tabular}{clrrrl}
\hline
$a$ & Name & $\Delta r$ & $\Delta K$ & $c(a)$ & Ecological interpretation \\
\hline
0 & Do Nothing                  & $0$     & $0$    & $0.00$  & Passive monitoring; natural baseline \\
1 & Aggressive Harvest          & $-0.02$ & $0$    & $-0.40$ & Heavy culling; revenue offsets cost \\
2 & Predator / Disease Control  & $+0.01$ & $0$    & $+0.20$ & Reduces mortality; raises $r_{\mathrm{eff}}$ \\
3 & Intensive Restoration       & $0$     & $+200$ & $+0.30$ & Major habitat works; raises $K_{\mathrm{eff}}$ \\
4 & Flagship Conservation       & $+0.02$ & $+200$ & $+0.60$ & Breeding + intensive restoration combined \\
\hline
\end{tabular}
\end{center}

\subsubsection*{2.2 Full Configuration (10 Actions)}

The full set adds intermediate-strength actions and additional combined interventions, making the optimal policy non-trivial under cost constraints.

\begin{center}
\resizebox{1.085\linewidth}{!}{%
\begin{tabular}{clrrrl}
\hline
$a$ & Name & $\Delta r$ & $\Delta K$ & $c(a)$ & Ecological interpretation \\
\hline
0 & Do Nothing                          & $0$     & $0$    & $0.00$  & Passive monitoring; natural baseline \\
1 & Sustainable Harvest                 & $-0.01$ & $0$    & $-0.20$ & Light culling / quota fishing; small revenue \\
2 & Aggressive Harvest                  & $-0.02$ & $0$    & $-0.40$ & Heavy culling; larger revenue, larger penalty \\
3 & Predator / Disease Control          & $+0.01$ & $0$    & $+0.20$ & Predator removal / pathogen treatment \\
4 & Breeding / Recruitment Support      & $+0.02$ & $0$    & $+0.40$ & Captive breeding, supplementary feeding \\
5 & Moderate Restoration                & $0$     & $+100$ & $+0.15$ & Fencing, fire management, planting \\
6 & Intensive Restoration               & $0$     & $+200$ & $+0.30$ & Reserve expansion, revegetation \\
7 & Integrated Conservation (light)     & $+0.01$ & $+100$ & $+0.35$ & Predator control + moderate restoration \\
8 & Adaptive Conservation Trial         & $+0.01$ & $+200$ & $+0.50$ & Predator control + intensive restoration \\
9 & Flagship Conservation Programme     & $+0.02$ & $+200$ & $+0.60$ & Breeding + intensive restoration combined \\
\hline
\end{tabular}%
}
\end{center}

\phantomsection\addcontentsline{toc}{subsection}{3. Transition Dynamics}
\subsection*{3. Transition Dynamics}
The continuous next state is computed via the Ricker map~\cite{ricker1954} under the effective parameters:
\[
    s_{t+1} = s_t \cdot \exp\!\left( r_{\mathrm{eff}}(a_t) \cdot \left(1 - \frac{s_t}{K_{\mathrm{eff}}(a_t)} \right) \right).
\]
Process and observation noise are set to zero in the current configuration ($\eta_t = 0$, $\epsilon_t = 0$); all stochasticity perceived by the agent arises from state aliasing introduced by discretisation and from the hidden episode-level draw of $r_{\mathrm{base}}$.

\phantomsection\addcontentsline{toc}{subsection}{4. Reward Function}
\subsection*{4. Reward Function}
The reward depends only on the current state and current action (predefined, not learned):
\[
    R(x_t, a_t) = \alpha \cdot \frac{x_t}{x_{\max}} - c(a_t), \qquad \alpha = 1, \quad x_{\max} = 100.
\]
The first term rewards keeping the discrete abundance bin high; the second term internalises the cost (or revenue, if negative) of the management action. This makes the optimal policy genuinely non-trivial even with $r_{\mathrm{base}}$ known: aggressive interventions (e.g.\ $a_9$) yield large ecological gains but at large cost, so the agent must weigh the marginal ecological return against the marginal monetary expenditure.

\phantomsection\addcontentsline{toc}{subsection}{5. State Space Discretization}
\subsection*{5. State Space Discretization}\begin{itemize}
    \item \textbf{Maximum Continuous Abundance:} $1000$.
    \item \textbf{Discrete States ($\mathcal{X}$):} $100$ intervals of width $10$, indexed $x \in \{0, 1, \dots, 99\}$.
    \item \textbf{Discretisation rule for logging:} $x_t = \lfloor s_t / 10 \rfloor$ (clipped to $\mathcal{X}$).
\end{itemize}

\phantomsection\addcontentsline{toc}{subsection}{6. The Trajectory Generation Loop}
\subsection*{6. The Trajectory Generation Loop}
The environment generates the offline dataset $\mathcal{D} = \{(x_t, a_t, R_t, x_{t+1})\}$ via the following loop:
\begin{enumerate}
    \item Sample a hidden growth rate $r_{\mathrm{base}} \sim \mathcal{U}(0.95, 1.00)$ for the episode.
    \item Sample a continuous starting state $s_0 \sim \mathcal{U}(1, 999)$.
    \item At each timestep $t$:
    \begin{enumerate}
        \item Compute the discrete state $x_t = \lfloor s_t / 10 \rfloor$.
        \item Select an action $a_t \in \mathcal{A}$ via a uniform-random (or suboptimal) data-collection policy.
        \item Look up $\Delta r(a_t)$ and $\Delta K(a_t)$ from Table 2.1 (smaller) or Table 2.2 (full); set $r_{\mathrm{eff}} = r_{\mathrm{base}} + \Delta r(a_t)$ and $K_{\mathrm{eff}} = K_{\mathrm{base}} + \Delta K(a_t)$.
        \item Compute the next continuous state $s_{t+1}$ via the Ricker equation above.
        \item Compute the immediate reward $R_t = R(x_t, a_t)$.
        \item Compute the next discrete state $x_{t+1} = \lfloor s_{t+1} / 10 \rfloor$.
        \item Log the tuple $(x_t, a_t, R_t, x_{t+1})$ to $\mathcal{D}$.
    \end{enumerate}
    \item The continuous states $\{s_t\}$ and the hidden $r_{\mathrm{base}}$ are discarded from the logged dataset; the agent only ever observes the discrete tuples.
\end{enumerate}

\newpage

\phantomsection\addcontentsline{toc}{section}{Part 2: The Black-Box Ensemble Architecture}
\section*{Part 2: The Black-Box Ensemble Architecture}
For this setting, a \textbf{Categorical Probabilistic Ensemble} is implemented, inspired by the MOPO architecture~\cite{yu2020mopo} and modified for discrete states. Since the states are binned, a Categorical Softmax Output is utilized.

\phantomsection\addcontentsline{toc}{subsection}{The Model: Bootstrapped MLP Ensemble with Embeddings}
\subsection*{The Model: Bootstrapped MLP Ensemble with Embeddings}
The architecture consists of $N = 5$ identical, independent neural networks. The configuration for one of those networks (using PyTorch terminology) is as follows:

\paragraph{1. The Input Layer (Embeddings)}
Because the states and actions are discrete, embedding layers are used rather than raw integers to avoid incorrect assumptions regarding the scale of the state indices.
\begin{quote}
    \texttt{StateEmbedding = nn.Embedding(num\_embeddings=100, embedding\_dim=32)}\\
    \texttt{ActionEmbedding = nn.Embedding(num\_embeddings=5, embedding\_dim=16)}
\end{quote}
These vectors are concatenated to produce a flat input vector of size 48. (The
\texttt{num\_embeddings} for the action table is set to the size of the active
action space: $5$ for the smaller configuration benchmarked in Part~3, or $10$ for
the full configuration of Table~2.2.)

\paragraph{2. State Aliasing Buffer (Frame Stacking)}
To provide velocity context and mitigate POMDP observation noise from discretization, the last two time-steps are provided as input:
\begin{quote}
    \texttt{Network Input: [StateEmbed(t-1), ActionEmbed(t-1), StateEmbed(t), ActionEmbed(t)]}
\end{quote}

\paragraph{3. The Hidden Layers}
A standard Multi-Layer Perceptron (MLP) is used for data efficiency.
\begin{quote}
    \texttt{Layer 1: Linear(96, 256) + ReLU()}\\
    \texttt{Layer 2: Linear(256, 256) + ReLU()}\\
    \texttt{Layer 3: Linear(256, 256) + ReLU()}
\end{quote}

\paragraph{4. The Output Layer (Categorical Distribution)}
\begin{quote}
    \texttt{Output Layer: Linear(256, 100)}
\end{quote}
The output is a vector of 100 logits. Applying a Softmax function yields a probability distribution over all 100 possible next states.

\phantomsection\addcontentsline{toc}{subsection}{The Training Process}
\subsection*{The Training Process}
To achieve the ensemble effect required for uncertainty detection:
\begin{enumerate}
    \item Five networks are initialized with different random weights.
    \item The offline dataset $\mathcal{D}$ is shuffled.
    \item Each network is trained on a different 80\% split of the dataset.
    \item \texttt{CrossEntropyLoss} is used to treat next-state prediction as a 100-class classification problem.
\end{enumerate}

\phantomsection\addcontentsline{toc}{subsection}{Rationale for the Architecture}
\subsection*{Rationale for the Architecture}
\begin{itemize}
    \item \textbf{Model-Agnostic:} The network is a standard MLP and requires no prior knowledge of ecological parameters.
    \item \textbf{Data Efficiency:} Through the use of embeddings and a classification framework, the network learns the underlying growth dynamics efficiently from moderate data.
    \item \textbf{Uncertainty Quantification:} If the ensemble is queried regarding an infrequently visited state, the five networks will output divergent distributions. The variance across these distributions is calculated to trigger the pessimism penalty.
\end{itemize}

\newpage

\phantomsection\addcontentsline{toc}{section}{Part 3: Experiment Setting (Selected-Method Benchmark)}
\section*{Part 3: Experiment Setting (Selected-Method Benchmark)}
To validate the pipeline end to end and to position the general model-based
offline RL approach against the ecology literature, a short, single-seed
benchmark was run. The goal was \emph{not} a tuned scientific comparison but a
framework demonstration: confirming that the shared environment interface, the
shared dataset wiring, the reward contract, the runtime logging, and the GPU
metadata capture all work together, and obtaining a first read on relative
behaviour. All hyperparameters are handoff defaults; all numbers below come from a
single seed (42) at demo scale.

\phantomsection\addcontentsline{toc}{subsection}{1. Methods Compared}
\subsection*{1. Methods Compared}
Four methods were benchmarked, deliberately a subset of the full implemented
ladder. Two are \emph{general} model-based offline RL methods that learn the
dynamics from data and assume no population-model family; two are
\emph{ecology-specific} mechanistic baselines that assume the Ricker family.

\paragraph{General model-based offline RL.}
\begin{itemize}
    \item \textbf{MOPO}~\cite{yu2020mopo} --- the Categorical Probabilistic Ensemble of Part~2
    ($N=5$ networks, frame-stack $2$, bootstrap ratio $0.8$, $20$ training epochs)
    combined with a \emph{pessimistic} receding-horizon planner. For each
    candidate action the planner rolls the ensemble forward and penalises the
    discounted return by the ensemble disagreement,
    \[
        \hat{r}(s,a) = \mathbb{E}[\text{reward}] - \lambda \cdot \mathrm{Var}_i\!\big(\mathbb{E}_i[x']\big),
        \qquad \lambda = 1.0,
    \]
    with planning horizon $5$ and $20$ Monte-Carlo rollouts.
    \item \textbf{RefPlan} (Reflect-then-Plan)~\cite{jeong2025refplan} --- a finite-state model-based
    planner (ensemble size $10$, planning horizon $4$, $64$ candidate action
    sequences drawn from a behaviour prior, $4$ latent samples, $\kappa = 5.0$,
    uncertainty penalty $0.10$). It plans over a learned latent posterior and is
    explicitly uncertainty-aware.
\end{itemize}

\paragraph{Ecology-specific mechanistic baselines (assume the Ricker family).}
\begin{itemize}
    \item \textbf{PLUS} (BioConserv'18)~\cite{memarzadeh2018} --- a finite hidden-model Bayesian
    baseline. It maintains a posterior over $21$ candidate values of $\rbase$
    spanning the prior, builds one transition matrix per candidate (transition
    grid $2000$, likelihood floor $10^{-12}$), solves each by exact value
    iteration ($\gamma = 0.95$), and acts by posterior-averaged Q-values. It does
    \emph{no} neural training.
    \item \textbf{MOOR-Ricker} (ExpertSys'23)~\cite{ju2023moor} --- a mechanistic offline RL
    baseline that \emph{learns} the growth parameters $(r, K)$ from the offline
    transitions by least squares ($10$ L-BFGS restarts, $10$ iterations each),
    discretises the learned Ricker model (transition grid $1000$), and plans by
    value iteration ($\gamma = 0.95$).
\end{itemize}

\phantomsection\addcontentsline{toc}{subsection}{2. Data, Evaluation, and Fairness}
\subsection*{2. Data, Evaluation, and Fairness}
A single offline dataset of $25{,}000$ transitions was generated per setting under
a uniform-random collection policy with seed $42$. MOPO, RefPlan, and MOOR-Ricker
consume the \emph{same} saved dataset for a given setting; PLUS does not train from
data but is evaluated on the same environment, seeds, action set, episode count,
horizon, and reward contract. All methods are evaluated on $50$ episodes of
horizon $50$ with discount $\gamma = 0.95$, using an identical sequence of
per-episode environment seeds ($42, \dots, 91$), so every method sees the same
evaluation episodes. The benchmark uses the smaller \textbf{5-action}
configuration of Table~2.1.

\phantomsection\addcontentsline{toc}{subsection}{3. The Two Settings and the Mismatch Mechanism}
\subsection*{3. The Two Settings and the Mismatch Mechanism}
\begin{description}
    \item[Primary (Ricker).] The true environment, the data, and the evaluation
    are all Ricker. PLUS uses the true Ricker transition (conditioned on each
    candidate $\rbase$) and MOOR-Ricker is well-specified.
    \item[Stress (Schaefer mismatch).] The true environment, the data, and the
    evaluation are all Schaefer/logistic~\cite{schaefer1954},
    \[
        s_{t+1} = s_t + \reff(a_t)\, s_t \left(1 - \frac{s_t}{\Keff(a_t)}\right),
    \]
    with identical states, actions, reward, episode length, and dataset interface.
    MOPO and RefPlan run \emph{unchanged} and learn from the Schaefer data. The
    mechanistic baselines are deliberately kept on their Ricker assumption: PLUS
    sets \texttt{use\_env\_transition=false} (its candidate matrices are built from
    the internal Ricker formula rather than the true Schaefer transition), and
    MOOR-Ricker fits a Ricker model (\texttt{dynamics\_model=ricker}) to the
    Schaefer-generated data.
\end{description}
The intended hypothesis is that the model-agnostic methods, which never assume a
population-model family, should remain competitive --- or pull ahead --- when the
assumed ecological model is wrong.

\phantomsection\addcontentsline{toc}{subsection}{4. Hardware and Scheduling}
\subsection*{4. Hardware and Scheduling}
Runs were submitted to Slurm as a four-task method array (task 0 = MOPO,
1 = RefPlan, 2 = PLUS, 3 = MOOR-Ricker). The Ricker methods executed on NVIDIA
L40S nodes; in the Schaefer setting the three light methods ran on a Tesla~T4.
MOPO and RefPlan used \texttt{device=cuda}; PLUS and MOOR-Ricker used
\texttt{device=cpu} within GPU jobs. A wrapper recorded wall-clock runtime, host,
GPU name, and a GPU-memory sample per method. The Schaefer MOPO run first hit the
earlier $4$\,h Slurm wall on the slower T4 (state \texttt{TIMEOUT}), then was
re-run with the \emph{identical} full-budget configuration on an L40S and
completed; its result is therefore full-budget and directly comparable.

\newpage

\phantomsection\addcontentsline{toc}{section}{Part 4: Results}
\section*{Part 4: Results}
Tables are sorted by mean cumulative reward over the $50$ evaluation episodes; the
reward cell also reports the episode-level standard deviation. ``Catastrophic
rate'' is the fraction of episodes whose minimum abundance bin is $\le 5$.

\phantomsection\addcontentsline{toc}{subsection}{1. Primary Setting (Ricker)}
\subsection*{1. Primary Setting (Ricker)}
\begin{center}
\begin{tabular}{lcccccc}
\toprule
Method & Cum.\ reward & $[\min,\max]$ & Final & Min & Cat. & Runtime \\
       & (mean\,$\pm$\,std) &        & abund. & abund. & rate & (s) \\
\midrule
MOOR-Ricker & $16.436 \pm 0.423$ & $[14.69, 16.87]$ & 49.5 & 38.2 & 0.06 & 158 \\
PLUS        & $16.436 \pm 0.423$ & $[14.69, 16.87]$ & 49.5 & 38.2 & 0.06 & 253 \\
MOPO        & $16.013 \pm 0.680$ & $[13.66, 16.69]$ & 49.5 & 37.8 & 0.06 & 10206 \\
RefPlan     & $14.947 \pm 0.763$ & $[12.61, 16.20]$ & 49.4 & 38.2 & 0.06 & 242 \\
\bottomrule
\end{tabular}
\end{center}
\noindent\textit{Primary Ricker setting, 50 evaluation episodes, seed 42; all runs
on an NVIDIA L40S.}

\phantomsection\addcontentsline{toc}{subsection}{2. Stress Setting (Schaefer Mismatch)}
\subsection*{2. Stress Setting (Schaefer Mismatch)}
\begin{center}
\begin{tabular}{lcccccc}
\toprule
Method & Cum.\ reward & $[\min,\max]$ & Final & Min & Cat. & Runtime \\
       & (mean\,$\pm$\,std) &        & abund. & abund. & rate & (s) \\
\midrule
MOOR-Ricker & $16.323 \pm 0.480$ & $[14.24, 16.66]$ & 50.0 & 32.7 & 0.08 & 24 \\
PLUS        & $16.323 \pm 0.480$ & $[14.24, 16.66]$ & 50.0 & 32.7 & 0.08 & 127 \\
MOPO        & $15.595 \pm 1.088$ & $[10.99, 16.62]$ & 50.0 & 34.9 & 0.08 & 10365 \\
RefPlan     & $14.688 \pm 0.835$ & $[11.86, 15.90]$ & 49.7 & 32.2 & 0.08 & 324 \\
\bottomrule
\end{tabular}
\end{center}
\noindent\textit{Schaefer model-mismatch stress setting, 50 evaluation episodes,
seed 42. PLUS, MOOR-Ricker, and RefPlan ran on a Tesla~T4; MOPO is the full-budget
re-run on an NVIDIA L40S after the original T4 task timed out at the 4\,h wall.
PLUS and MOOR-Ricker retain a Ricker-family assumption against a Schaefer truth.}

\phantomsection\addcontentsline{toc}{subsection}{3. Epistemic Uncertainty and Policy Entropy}
\subsection*{3. Epistemic Uncertainty and Policy Entropy}
These two auxiliary signals separate the methods far more than reward does.

\begin{center}
\begin{tabular}{lcccc}
\toprule
 & \multicolumn{2}{c}{Ricker (primary)} & \multicolumn{2}{c}{Schaefer (stress)} \\
\cmidrule(lr){2-3}\cmidrule(lr){4-5}
Method & uncertainty & entropy & uncertainty & entropy \\
\midrule
MOOR-Ricker & 0.0000 & 0.007 & 0.0000 & 0.014 \\
PLUS        & 0.0028 & 0.015 & 0.0039 & 0.024 \\
MOPO        & 0.0425 & 0.249 & 0.1468 & 0.422 \\
RefPlan     & 0.0668 & 0.370 & 0.2207 & 0.520 \\
\bottomrule
\end{tabular}
\end{center}
\noindent\textit{Mean logged epistemic uncertainty and mean policy entropy. Both
learned methods become much less certain under the Schaefer mismatch (RefPlan
$0.067 \to 0.221$, max $2.25 \to 6.46$; MOPO $0.043 \to 0.147$, max
$1.56 \to 5.02$), while the mechanistic baselines stay confident.}

\phantomsection\addcontentsline{toc}{subsection}{4. Compute and Runtime}
\subsection*{4. Compute and Runtime}
Runtime is dominated by MOPO's planning evaluation, not its model fitting (the
ensemble trains and checkpoints in $\sim$3\,min). Each decision step issues
\[
    \underbrace{5}_{\text{actions}} \times \underbrace{20}_{\text{rollouts}}
    \times \underbrace{\textstyle\sum_{h=0}^{4}\big(1 + 5\cdot\mathbf{1}[h<4]\big)}_{=25}
    = 2500
\]
single-sample ensemble \texttt{predict} calls, so a $50$-episode, $50$-step
evaluation issues $\approx 6.25 \times 10^{6}$ ensemble forward passes. One
\texttt{predict} call was measured at $\sim$28.5\,ms on CPU and inferred at
$\sim$1.9\,ms on the L40S, giving $\sim$2.8\,h for the Ricker MOPO evaluation
(\texttt{elapsed} $=10206$\,s). The remaining methods each finished in seconds to
minutes. GPU-memory accounting was unreliable for the original batch (node-level
sampler returned $0$\,MB for most methods); a per-process sampler added for the
Schaefer MOPO re-run reported a peak of $\sim$564\,MB.

\newpage

\phantomsection\addcontentsline{toc}{section}{Part 5: Analysis}
\section*{Part 5: Analysis}
\phantomsection\addcontentsline{toc}{subsection}{1. The Ecology Baselines Lead; the Learned Planners Pay an Approximation Tax}
\subsection*{1. The Ecology Baselines Lead; the Learned Planners Pay an Approximation Tax}
On the primary Ricker task PLUS and MOOR-Ricker achieve the highest reward, MOPO
is close behind, and RefPlan trails. This ordering is consistent with the nature
of the task: a \emph{well-specified, near-deterministic, exact-observation finite
MDP}. Given $\rbase$ the dynamics are deterministic, and the only stochasticity
the agent perceives is the narrow prior over $\rbase$ plus bin aliasing. PLUS
knows the Ricker form and the correct $K_{\mathrm{base}}$ and solves each candidate
MDP exactly; MOOR-Ricker recovers $(r, K)$ accurately by least squares and likewise
plans by exact value iteration. Both effectively \emph{solve} the MDP. MOPO and
RefPlan, by contrast, approximate the dynamics with a neural ensemble or latent
model and plan with short horizons and explicit pessimism; where the true model is
easy to identify, that approximation error and built-in conservatism cost reward
rather than protect against it. MOPO beating RefPlan is consistent with MOPO's
held-out dynamics accuracy ($0.79$) and its value-based pessimistic rollout
extracting more reward than RefPlan's shorter-horizon, more heavily penalised
latent planning.

\phantomsection\addcontentsline{toc}{subsection}{2. PLUS and MOOR-Ricker Are Numerically Identical}
\subsection*{2. PLUS and MOOR-Ricker Are Numerically Identical}
In \emph{both} settings PLUS and MOOR-Ricker produce identical per-episode
cumulative rewards (Ricker $16.436$; Schaefer $16.323$) and hence identical
abundance statistics. This is not a bug. With exact observations and a shared
reward contract, both converge to the same greedy policy: PLUS via a posterior over
$\rbase$ concentrated on the correct Ricker form, MOOR via a least-squares Ricker
fit. They differ only in their auxiliary uncertainty signals. In this
exact-observation regime the two ecology baselines are effectively redundant; they
may diverge under harder (partially observed or noisier) settings.

\phantomsection\addcontentsline{toc}{subsection}{3. Why the Mismatch Stress Did Not Expose the Intended Advantage}
\subsection*{3. Why the Mismatch Stress Did Not Expose the Intended Advantage}
The mismatch did \emph{not} favour the model-agnostic methods: the
Ricker-assuming baselines fell only $0.11$ ($16.44 \to 16.32$), RefPlan fell $0.26$
($14.95 \to 14.69$), and MOPO fell the \emph{most}, $0.42$ ($16.01 \to 15.60$);
the within-setting ordering is unchanged. The root cause is that the mismatch is
largely \emph{decision-irrelevant}:
\begin{enumerate}
    \item \textbf{Shared equilibrium and stability.} Both maps have fixed points at
    $s=0$ and $s=K$, and the same linearisation $f'(K) = 1 - r$ at carrying
    capacity, which is stable for $r \in (0,2)$. For $\rbase \in [0.95, 1.00]$ both
    families converge monotonically to the same $\Keff$, differing only in the
    low-abundance transient.
    \item \textbf{Family-invariant optimum.} Because the reward rewards sitting near
    high abundance and avoiding costly actions, and both families drive the
    population to the \emph{same} $\Keff$ set by the \emph{same} action metadata,
    the optimal action ranking is essentially identical under either assumption.
    The mismatch perturbs the path, not the decision.
    \item \textbf{The baselines are not very ``wrong''.} MOOR re-fits $(r,K)$ to the
    Schaefer data and lands near the true $K$; PLUS uses the correct
    $K_{\mathrm{base}}$ and $\rbase$ grid. Only the functional form differs, and the
    two forms agree at the operating point $K$.
    \item \textbf{Coarse discretisation and a narrow prior} ($100$ bins of width
    $10$; $\rbase \in [0.95, 1.0]$) further wash out the transient differences.
\end{enumerate}
The clearest evidence is that PLUS and MOOR-Ricker remain numerically identical
across both settings --- a direct demonstration that the model family does not
change the decision here.

\phantomsection\addcontentsline{toc}{subsection}{4. Under Mismatch, Uncertainty Rises --- Aiding Calibration but Costing Reward}
\subsection*{4. Under Mismatch, Uncertainty Rises --- Aiding Calibration but Costing Reward}
Both learned methods become markedly less certain when the truth becomes Schaefer
(RefPlan $0.067 \to 0.221$, MOPO $0.043 \to 0.147$), while the mechanistic
baselines stay confident. This has two opposite consequences. Positively, it is a
genuine \emph{calibration} signal: the model-agnostic methods \emph{detect} that
they are out of distribution, whereas PLUS and MOOR remain overconfident with a
wrong-but-adequate Ricker model. Negatively, it directly costs reward: MOPO
subtracts $\lambda \cdot \mathrm{Var}_i$ from every imagined step ($\lambda = 1$)
and RefPlan applies an explicit uncertainty penalty, so greater disagreement makes
both plan more conservatively. This explains why MOPO loses the most reward under
mismatch: the unfamiliar dynamics inflate ensemble disagreement, and pessimism
converts that disagreement into caution even though no catastrophe is present in
this benign regime. The mismatch here does not reward model-agnosticism --- it
\emph{penalises} the uncertainty-aware methods on reward while improving their
calibration.

\phantomsection\addcontentsline{toc}{subsection}{5. Safety Profile}
\subsection*{5. Safety Profile}
All methods keep the population healthy in both settings: mean abundance sits near
bin $49$--$50$ (the population equilibrates near $K = 500 \Rightarrow$ bin $50$),
and catastrophic episodes are rare ($0.06$ Ricker, $0.08$ Schaefer). The slightly
higher catastrophic rate and lower minimum abundance ($\sim$32 vs.\ $\sim$38)
under Schaefer reflect the logistic's slower recovery from low abundance, but no
method collapses. The task lives in a ``safe'' region where the abundance-maximising
and safety objectives are aligned --- a further reason the mismatch is not punishing.

\newpage

\phantomsection\addcontentsline{toc}{section}{Part 6: Potential Improvements}
\section*{Part 6: Potential Improvements}
The benchmark confirms the harness runs end to end but does not yet exercise the
regime where a model-agnostic learner should win. The improvements below target
that gap, plus the compute bottleneck and the statistical weakness of a single
seed.

\phantomsection\addcontentsline{toc}{subsection}{1. Designing Scenarios That Favour the Model-Agnostic Methods}
\subsection*{1. Designing Scenarios That Favour the Model-Agnostic Methods}
A family-agnostic learner (RefPlan, MOPO) can out-score a baseline that assumes one
or two known families only when the committed family assumption is both \emph{wrong}
and \emph{costly}. Concretely, a scenario rewards the agnostic learner when all four
of the following hold:
\begin{enumerate}
    \item \textbf{Out-of-family truth.} The true transition law lies outside every
    assumed family. If the baseline assumes $\{$Ricker, logistic$\}$, the truth must
    be a \emph{third} form (or a mixture / regime switch); any model inside the
    assumed set can simply be fitted. The Schaefer stress partly failed here because
    the logistic is itself a family one would naturally assume.
    \item \textbf{Decision-relevance.} The misspecification must change the
    \emph{optimal action}, not merely the transient path. Misspecification at the
    operating point is harmless if both models share the equilibrium and the reward
    depends only on it.
    \item \textbf{Learnability from data.} The flexible model is accurate only where
    the offline data covers the relevant $(x,a)$ cells. The behaviour policy must
    therefore \emph{visit} the decision-critical region (e.g.\ low abundance).
    \item \textbf{Asymmetric / catastrophic cost.} The advantage is largest when
    getting that region wrong is expensive (collapse, extinction) and the region is
    partly data-sparse, so the uncertainty penalty steers the agnostic learner away
    while a confident parametric extrapolation walks in.
\end{enumerate}
Conditions 1--2 make the family assumption \emph{wrong}; conditions 3--4 let the
flexible learner \emph{exploit} it.

\paragraph{True-environment families outside Ricker/logistic.}
Each keeps the same states, actions, reward, and hidden $\rbase$; only the update
map changes (a one-method subclass of the environment, exactly as the Schaefer
variant already is):
\begin{description}
    \item[Critical depensation / Allee~\cite{liermann2001, courchamp2008}.]
    $s_{t+1} = s_t \exp\!\big(\reff(1 - s_t/\Keff)(s_t/C - 1)\big)$ with a critical
    threshold $C$: growth is \emph{negative} below $C$, so the population collapses
    once pushed under it. Compensatory families (Ricker, logistic) have no such
    threshold, so a family-assuming agent harvests through $C$ expecting recovery.
    This is the canonical ``wrong model $\Rightarrow$ catastrophic over-harvest''
    case, and depensation is empirically documented in depleted fish
    stocks~\cite{liermann2001}.
    \item[$\theta$-logistic~\cite{gilpin1973theta}.]
    $s_{t+1} = s_t + \reff\, s_t\big(1 - (s_t/\Keff)^{\theta}\big)$ with
    $\theta \neq 1$ concentrates or spreads density dependence; neither Ricker nor
    logistic ($\theta = 1$) matches the curvature, and the optimal harvest level
    shifts with $\theta$.
    
\end{description}

\paragraph{Experimental knobs that widen the gap.}
Replace the pure random collection policy with one that also visits low abundance
(so the danger zone is in-distribution); sharpen the reward on collapse (an explicit
penalty or an absorbing extinction threshold, since the current contract barely
penalises low abundance at $6$--$8\%$ catastrophic rate); push toward the regime
where transient $\neq$ equilibrium (widen the hidden prior and/or raise $\reff$
toward the over-compensation range $r \gtrsim 2$); add heteroscedastic or
heavy-tailed process/observation noise where mistakes are costly; and lengthen the
evaluation horizon so threshold/transient errors compound.

\paragraph{A concrete protocol.}
A minimal high-leverage instance: true environment $=$ critical-depensation Ricker
with threshold $C \approx 0.3\,K$; baselines $=$ PLUS and MOOR-Ricker kept on the
compensatory Ricker family (optionally give PLUS a second candidate family,
logistic, to test the ``two known models'' case --- both still lack depensation);
MOPO and RefPlan unchanged; a collection policy that drives the population near and
below $C$ so the offline data contains collapses; and an added extinction penalty.
Hypothesis: the family-assuming baselines, lacking any representation of the
threshold, over-harvest into collapse on a non-trivial fraction of episodes,
whereas the agnostic learner --- having learned the negative-growth region and
penalising uncertain, data-sparse states --- keeps abundance above $C$ and attains
higher return at a lower catastrophic rate.

\phantomsection\addcontentsline{toc}{subsection}{2. Compute and Implementation}
\subsection*{2. Compute and Implementation}
MOPO's full-budget planning evaluation should either be given a longer wall on a
fast GPU, run at a disclosed reduced planning budget
($\texttt{num\_rollouts}=1$, $\texttt{horizon}=3$ cuts the cost $\sim$40$\times$ to
$\sim$1.5\,s/step), or --- for a real benchmark --- accelerated by vectorising the
\texttt{predict} call across ensemble members and rollouts, turning
$\sim$28\,ms/call into sub-millisecond and finishing full budget in minutes.
Caching repeated finite-state rollouts would further help short-turnaround
iteration.

\phantomsection\addcontentsline{toc}{subsection}{3. Statistical Rigour and Fairness}
\subsection*{3. Statistical Rigour and Fairness}
Firm conclusions require \emph{multiple seeds} and \emph{tuned hyperparameters};
the present single demo seed only indicates a direction. When constructing the
harder stress tests above, keep states, actions, reward, seeds, episode count, and
horizon identical across methods, and give the parametric baselines their
\emph{best} fit --- the win must come from family misspecification, not from
under-fitting the baselines. Ensure the deviation is learnable (the data must visit
the out-of-family region) while avoiding dynamics so noisy or data-starved that
\emph{every} method fails.

\newpage

\phantomsection
\addcontentsline{toc}{section}{References}
\begin{thebibliography}{99}

\bibitem{yu2020mopo}
T.~Yu, G.~Thomas, L.~Yu, S.~Ermon, J.~Y.~Zou, S.~Levine, C.~Finn, and T.~Ma.
\newblock MOPO: Model-Based Offline Policy Optimization.
\newblock In \emph{Advances in Neural Information Processing Systems (NeurIPS)}, 2020.
\newblock arXiv:2005.13239. \url{https://arxiv.org/abs/2005.13239}

\bibitem{kidambi2020morel}
R.~Kidambi, A.~Rajeswaran, P.~Netrapalli, and T.~Joachims.
\newblock MOReL: Model-Based Offline Reinforcement Learning.
\newblock In \emph{Advances in Neural Information Processing Systems (NeurIPS)}, 2020.
\newblock arXiv:2005.05951. \url{https://arxiv.org/abs/2005.05951}

\bibitem{jeong2025refplan}
J.~Jeong, X.~Wang, J.~Wang, S.~Sanner, and P.~Poupart.
\newblock Reflect-then-Plan: Offline Model-Based Planning through a Doubly Bayesian Lens.
\newblock In \emph{Proceedings of the 42nd International Conference on Machine Learning (ICML)}, 2025 (Spotlight).
\newblock arXiv:2506.06261. \url{https://arxiv.org/abs/2506.06261}

\bibitem{memarzadeh2018}
M.~Memarzadeh and C.~Boettiger.
\newblock Adaptive management of ecological systems under partial observability.
\newblock \emph{Biological Conservation}, 224:9--15, 2018.
\newblock doi:10.1016/j.biocon.2018.05.009.

\bibitem{chades2012momdp}
I.~Chad\`es, J.~Carwardine, T.~G.~Martin, S.~Nicol, R.~Sabbadin, and O.~Buffet.
\newblock MOMDPs: A Solution for Modelling Adaptive Management Problems.
\newblock In \emph{Proceedings of the AAAI Conference on Artificial Intelligence}, 2012.

\bibitem{ju2023moor}
Y.~Ju et al.
\newblock Model-based offline reinforcement learning for sustainable fishery management.
\newblock \emph{Expert Systems}, Wiley, 2023.
\newblock doi:10.1111/exsy.13324.

\bibitem{ricker1954}
W.~E.~Ricker.
\newblock Stock and recruitment.
\newblock \emph{Journal of the Fisheries Research Board of Canada}, 11(5):559--623, 1954.

\bibitem{schaefer1954}
M.~B.~Schaefer.
\newblock Some aspects of the dynamics of populations important to the management of the commercial marine fisheries.
\newblock \emph{Bulletin of the Inter-American Tropical Tuna Commission}, 1(2):23--56, 1954.

\bibitem{gilpin1973theta}
M.~E.~Gilpin and F.~J.~Ayala.
\newblock Global models of growth and competition.
\newblock \emph{Proceedings of the National Academy of Sciences (PNAS)}, 70(12):3590--3593, 1973.

\bibitem{liermann2001}
M.~Liermann and R.~Hilborn.
\newblock Depensation: evidence, models and implications.
\newblock \emph{Fish and Fisheries}, 2(1):33--58, 2001.

\bibitem{courchamp2008}
F.~Courchamp, L.~Berec, and J.~Gascoigne.
\newblock \emph{Allee Effects in Ecology and Conservation}.
\newblock Oxford University Press, 2008.

\bibitem{turchin1990}
P.~Turchin.
\newblock Rarity of density dependence or population regulation with lags?
\newblock \emph{Nature}, 344:660--663, 1990.

\end{thebibliography}

\end{document}